% !TEX program = pdflatex
\documentclass[12pt,a4paper]{letter}

\usepackage[T1]{fontenc}
\usepackage[margin=1in]{geometry}
\usepackage{hyperref}
\usepackage{setspace}
\usepackage{parskip}

\hypersetup{
  colorlinks=true,
  urlcolor=blue
}

\onehalfspacing

\signature{Professor {[FULL NAME]}\\
Professor of Computer Science\\
{[Department / Faculty]}\\
University of Ibadan, Nigeria}

\address{{[Department of Computer Science]}\\
University of Ibadan\\
Ibadan, Oyo State\\
Nigeria}

\date{6 March 2026}

\begin{document}

\begin{letter}{UK Visas \& Immigration\\
Global Talent Visa --- Tech Nation (Digital Technology)\\
United Kingdom}

\opening{Dear Sir/Madam,}

I write to endorse Odeyemi Olusegun Israel for the UK Global Talent
Visa in Digital Technology under the Exceptional Promise pathway.
My endorsement is based on the combination of \textbf{original research,
systems architecture, and practical engineering execution} that he has
demonstrated across his work.

I am a Professor of Computer Science at the University of Ibadan, where my
research includes {[relevant areas, e.g., machine learning,
distributed systems, software architecture, cybersecurity]}.  My
assessment is based on an \textbf{independent academic review} of
Olusegun's published preprints, technical evidence portfolio,
architectural materials, and open technical work.  Because his papers
have been released first as preprints rather than through formal peer
review, I consider this kind of external expert evaluation especially
important.  Having examined the substance of the work directly, I am
satisfied that it reflects genuine originality, technical depth, and an
unusual ability to translate research ideas into a live, multi-layered
technical system.

\textbf{AMTTP as a substantial architecture and engineering achievement.}

The strongest practical demonstration of Olusegun's capability is the
Anti-Money Laundering Transaction Transfer Protocol (AMTTP), which he
has designed and built as a full-stack compliance platform for digital
asset transactions.  This is not a superficial prototype.  It combines a
cross-platform Flutter application, a Next.js enterprise dashboard, a
FastAPI compliance orchestration layer, multiple specialist services for
machine-learning risk scoring, graph analysis, sanctions evaluation,
monitoring, georisk, data integrity, and explainability, together with a
substantial Solidity smart-contract layer for on-chain execution.

From a computer science and systems-design perspective, this is a
well-structured architecture.  The platform separates user interaction,
service orchestration, analytical processing, policy evaluation, and
blockchain enforcement in a way that reflects mature engineering
judgment.  I was particularly struck by the fact that AMTTP addresses
compliance in the \textit{pre-settlement} window rather than treating it
as a purely post-transaction monitoring problem.  That is a materially
harder technical challenge, because it requires risk assessment,
graph-based reasoning, and policy enforcement to operate inside real
transaction timing constraints.  In my view, this is a serious piece of
technical architecture rather than a conceptual claim.

His engineering decisions within that architecture are equally
important.  The evidence shows thoughtful treatment of performance,
deployability, and maintainability.  In particular, his use of
knowledge-distilled student models demonstrates an understanding that
machine-learning quality must be balanced against latency and operating
cost if a compliance system is to be usable in practice.  Likewise, his
integration of zero-knowledge compliance mechanisms into the broader
platform addresses a difficult systems problem at the intersection of
privacy, regulation, and distributed computing.

\textbf{Blind Spot Decomposition Theory (BSDT) as the core research
contribution.}

Alongside the architecture, the most important academic element of
Olusegun's profile is his development of Blind Spot Decomposition Theory
(BSDT).  This work is significant because it does not merely seek to
improve predictive accuracy in the aggregate.  Instead, it addresses the
harder and more practically important question of \textit{why}
well-performing models still miss the specific anomalies that matter most
in deployment.

BSDT provides a structured framework for analysing hidden failure modes
by decomposing missed detections into four components: Camouflage,
Feature Gap, Activity Anomaly, and Temporal Novelty.  That decomposition
is valuable because it turns model failure from an opaque residual into
something that can be analysed systematically.  In high-stakes domains
such as fraud detection and compliance, this matters greatly: the real
problem is often not average performance, but the existence of rare,
high-impact blind spots that create outsized operational risk.

The work is also technically credible in its validation.  As I
understand the evidence, the framework was evaluated under a strict
experimental design intended to avoid leakage and spurious improvement,
and it demonstrated large recall improvements without requiring the base
model itself to be retrained.  The reported correction performance---up
to \textbf{84.4 percentage points} improvement in recall in validation
contexts---is substantial.  More importantly, the work shows originality
of thought.  It is attempting to formalise the latent structure of
machine-learning failure, not simply to adjust thresholds or repackage
existing techniques.

What strengthens BSDT further is that it is not isolated from the rest
of his engineering work.  In Olusegun's case, the research informs a
broader system-building programme.  The same analytical mindset visible
in BSDT is reflected in how AMTTP approaches risk understanding,
compliance decisions, and explainability.  This connection between
research insight and technical implementation is one of the clearest
indicators that he is capable of moving from theory to real digital
infrastructure.

\textbf{My assessment.}

If I compare the intellectual ambition and technical execution of this
body of work to the level I normally see from early-career researchers
and engineers, Olusegun stands out very clearly.  He has not limited
himself to writing papers, nor has he confined himself to routine
software implementation.  He has combined original analytical thinking,
serious systems architecture, machine-learning engineering, and
privacy-aware blockchain infrastructure in a way that is unusual for a
single individual working independently.

In particular, I want to emphasize that my endorsement is not based on
informal familiarity alone, but on a direct evaluation of the technical
record itself.  Where work is presented through preprints, repositories,
and demonstrable system evidence rather than conventional journal
publication, the central question is whether an experienced academic can
independently examine that material and conclude that it is credible,
original, and significant.  In Olusegun's case, my answer to that
question is unequivocally yes.

The United Kingdom has strong interests in trustworthy AI, financial
technology, cybersecurity, and digitally enforceable compliance
infrastructure.  Olusegun's work aligns directly with those priorities.
I therefore consider him exceptionally well suited to endorsement under
the Exceptional Promise route, and I support his application without
reservation.

I endorse his application for the Global Talent Visa without
reservation.

\closing{Yours faithfully,}

\end{letter}

\end{document}
